\chapter{Analytické SQL dotazy a~indexy}
\label{chap:appendix:sql}

Tato příloha obsahuje zdrojový kód analytických \SQL dotazů použitých při evaluaci v~\customref{Sekci}{sec:evaluace:sql} a~definice databázových indexů pro jednotlivé tabulky.

% ─────────────────────────────────────────────────────────────────────────────
\section{Analytické dotazy}
\label{sec:appendix:sql:queries}

\subsection*{S1 — Počet \PR ve stavu (štítek) k~zadanému datu}
\label{lst:appendix:sql:s1}

Dotaz pomocí dvou \CTE s~\code{DISTINCT ON} určí, kolik \PR je v~daném stavu ke zvolenému datu. První \CTE \code{latest\_labels} zachová pro každý \PR nejnovější labelovou událost pro cílový štítek \code{S-*}, druhá \CTE \code{latest\_state} nejnovější stavovou změnu (uzavření, sloučení nebo znovuotevření). Výsledný počet zahrnuje pouze \PR, jejichž poslední labelová událost je \code{ADDED} a~které jsou ke zvolenému datu stále otevřené.

\begin{longlisting}
  \caption{S1 --- Počet \PR ve stavu (štítek) k~zadanému datu}
  \label{lst:appendix:sql:s1:code}
  \inputminted{sql}{./code/sql/q1_pr_count_by_label.sql}
\end{longlisting}

\subsection*{S2 — Kumulativní počet sloučených \PR k~zadanému datu}
\label{lst:appendix:sql:s2}

Na rozdíl od S1 nevyžaduje tento dotaz sledování historických přechodů stavů. Stav „merged" je trvalý — jakmile je \PR sloučen, nelze jej vrátit do předchozího stavu. Stačí proto spočítat počet unikátních \PR, u~nichž existuje záznam události \code{merged} v~tabulce \code{issue\_event\_history} s~časovým razítkem před zadaným datem. Výsledkem je \code{COUNT(DISTINCT issue)} s~indexovatelným filtrem \code{event = 'merged'} a~časovým řezem.

\begin{longlisting}
  \caption[S2 --- Kumulativní počet sloučených \PR]{S2 --- Kumulativní počet sloučených \PR k~zadanému datu}
  \label{lst:appendix:sql:s2:code}
  \inputminted{sql}{./code/sql/q2_merged_pr_count.sql}
\end{longlisting}

\subsection*{S3 — Soubory upravené přispěvateli týmu v~zadaném období}
\label{lst:appendix:sql:s3}

Dotaz agreguje soubory upravené přispěvateli daného týmu (\code{contributors\_teams}) v~zadaném časovém rozsahu. Výsledek je seřazen sestupně podle počtu úprav.

\begin{longlisting}
  \caption{S3 --- Soubory upravené přispěvateli týmu}
  \label{lst:appendix:sql:s3:code}
  \inputminted{sql}{./code/sql/q3_files_by_team.sql}
\end{longlisting}

\subsection*{S4 — Denní vývoj počtu \PR v~daném stavu v~časovém rozsahu}
\label{lst:appendix:sql:s4}

Nejsložitější z~testovaných dotazů využívá delta-event přístup. \CTE \code{all\_transitions} sestaví pomocí \code{UNION ALL} seznam všech relevantních událostí (vznik \PR, uzavření, sloučení, znovuotevření). \CTE \code{ordered\_transitions} přiřadí každé události čas následující události přes okenní funkci \code{LEAD()}, čímž vzniknou neprotínající se intervaly otevřenosti (\code{open\_periods}). Stejným způsobem jsou konstruovány intervaly aktivního štítku (\code{label\_active\_periods}). Průnik obou typů intervalů operátorem \code{\&\&} a~\code{*} dá intervaly, kdy je \PR v~cílovém stavu (\code{in\_state\_periods}). Tyto intervaly jsou rozloženy na delta-události (+1 na začátku, $-1$ na konci), agregovány po dnech a~pak je nad nimi spočítán kumulativní součet. Složitost je $O(\text{events} + \text{days})$ místo naivního $O(\text{days} \times \text{events})$.

\begin{longlisting}
  \caption{S4 --- Denní vývoj počtu \PR v~daném stavu}
  \label{lst:appendix:sql:s4:code}
  \inputminted{sql}{./code/sql/q4_pr_state_over_time.sql}
\end{longlisting}

% ─────────────────────────────────────────────────────────────────────────────
\section{Databázové indexy}
\label{sec:appendix:sql:indexes}

Níže jsou uvedeny definice všech databázových indexů. Indexy jsou rozděleny podle tabulky, ke které náleží.

\subsection*{Tabulka \code{issues}}
\begin{longlisting}
  \caption[Indexy tabulky issues]{Indexy tabulky \code{issues}}
  \label{lst:appendix:sql:idx_issues}
  \inputminted{sql}{./code/sql/issues_indexes.sql}
\end{longlisting}

\subsection*{Tabulka \code{issue\_event\_history}}
\begin{longlisting}
  \caption[Indexy tabulky issue\_event\_history]{Indexy tabulky \code{issue\_event\_history}}
  \label{lst:appendix:sql:idx_event_history}
  \inputminted{sql}{./code/sql/issue_event_history_indexes.sql}
\end{longlisting}

\subsection*{Tabulka \code{issue\_labels\_history}}
\begin{longlisting}
  \caption[Indexy tabulky issue\_labels\_history]{Indexy tabulky \code{issue\_labels\_history}}
  \label{lst:appendix:sql:idx_labels_history}
  \inputminted{sql}{./code/sql/issue_labels_history_indexes.sql}
\end{longlisting}

\subsection*{Tabulka \code{file\_activity}}
\begin{longlisting}
  \caption[Indexy tabulky file\_activity]{Indexy tabulky \code{file\_activity}}
  \label{lst:appendix:sql:idx_file_activity}
  \inputminted{sql}{./code/sql/file_activity_indexes.sql}
\end{longlisting}

\subsection*{Tabulka \code{contributors}}
\begin{longlisting}
  \caption[Indexy tabulky contributors]{Indexy tabulky \code{contributors}}
  \label{lst:appendix:sql:idx_contributors}
  \inputminted{sql}{./code/sql/contributors_indexes.sql}
\end{longlisting}

\subsection*{Tabulka \code{contributors\_teams}}
\begin{longlisting}
  \caption[Indexy tabulky contributors\_teams]{Indexy tabulky \code{contributors\_teams}}
  \label{lst:appendix:sql:idx_contributors_teams}
  \inputminted{sql}{./code/sql/contributors_teams_indexes.sql}
\end{longlisting}

% ─────────────────────────────────────────────────────────────────────────────
\section{Výstupy EXPLAIN ANALYZE}
\label{sec:appendix:sql:explain}

Níže jsou uvedeny výstupy příkazu \code{EXPLAIN (ANALYZE, BUFFERS)} pro každý testovaný dotaz, vždy bez indexů a~poté se zapnutými indexy. Měření proběhlo na databázi repozitáře \texttt{rust-lang/rust} s~přibližně 200\,000 \PR, 564\,000 záznamy v~\code{issue\_event\_history}, 355\,000 v~\code{issue\_labels\_history} a~525\,000 v~\code{file\_activity}.

% ── S1 ─────────────────────────────────────────────────────────────────────
\subsection*{S1 --- bez indexů}

Bez indexů provádí plánovač dva sekvenční průchody (\code{Seq Scan}) nad celými tabulkami \\ \code{issue\_labels\_history} a~\code{issue\_event\_history}. Výsledky jsou následně řazeny na disku operací \code{Sort (external merge Disk)}, protože data se nevejdou do operační paměti. Odhadované náklady dosahují 22\,258 při celkové době vykonávání 73,5\,ms.

\begin{longlisting}
  \caption{S1: EXPLAIN ANALYZE bez indexů (Execution Time: 73,5 ms)}
  \label{lst:appendix:explain:s1:no_idx}
  \inputminted[fontsize=\footnotesize]{text}{./code/sql/explain_q1_no_idx.txt}
\end{longlisting}

\subsection*{S1 --- s~indexy}

Hlavní změny oproti variantě bez indexů: oba \code{Seq Scan} + \code{Sort (external merge Disk)} nahrazeny dvěma \code{Index Only Scan}, řazení probíhá přímo po indexu.

\begin{longlisting}
  \caption{S1: EXPLAIN ANALYZE s~indexy (Execution Time: 20,0 ms)}
  \label{lst:appendix:explain:s1:idx}
  \inputminted[fontsize=\footnotesize]{text}{./code/sql/explain_q1_idx.txt}
\end{longlisting}

% ── S2 ─────────────────────────────────────────────────────────────────────
\subsection*{S2 --- bez indexů}

Bez indexů plánovač zvolí \code{Parallel Seq Scan} přes celou tabulku \code{issue\_event\_history} a~záznamy filtruje průběžně. Paralelizace sice snižuje elapsed time, avšak všechny záznamy jsou čteny bez ohledu na hodnotu sloupce \code{event}. Odhadované náklady jsou 9\,598 při době vykonávání 18,8\,ms.

\begin{longlisting}
  \caption{S2: EXPLAIN ANALYZE bez indexů (Execution Time: 18,8 ms)}
  \label{lst:appendix:explain:s2:no_idx}
  \inputminted[fontsize=\footnotesize]{text}{./code/sql/explain_q2_no_idx.txt}
\end{longlisting}

\subsection*{S2 --- s~indexy}

\code{Parallel Seq Scan} nahrazen přímým \code{Index Only Scan using idx\_pr\_event\_hist\_merged}. Paralelní scan zaniká, plán se zjednodušuje na jediný uzel.

\begin{longlisting}
  \caption{S2: EXPLAIN ANALYZE s~indexy (Execution Time: 2,7 ms)}
  \label{lst:appendix:explain:s2:idx}
  \inputminted[fontsize=\footnotesize]{text}{./code/sql/explain_q2_idx.txt}
\end{longlisting}

% ── S3 ─────────────────────────────────────────────────────────────────────
\subsection*{S3 --- bez indexů}

Bez indexů provede plánovač \code{Parallel Seq Scan} přes celou tabulku \code{file\_activity} a~záznamy nesouvisející s~daným týmem jsou čteny a~až poté filtrovány. Tabulka \code{contributors\_teams} je procházena sekvenčně a~výsledky jsou spojeny operací \code{Hash Join}. Odhadované náklady jsou 16\,770 při době vykonávání 50,6\,ms.

\begin{longlisting}
  \caption{S3: EXPLAIN ANALYZE bez indexů (Execution Time: 50,6 ms)}
  \label{lst:appendix:explain:s3:no_idx}
  \inputminted[fontsize=\footnotesize]{text}{./code/sql/explain_q3_no_idx.txt}
\end{longlisting}

\subsection*{S3 --- s~indexy}

\code{Parallel Seq Scan} na \code{file\_activity} nahrazen \code{Nested Loop} s~dvěma \code{Index Only Scan}: nejprve jsou načteni členové týmu z~primárního klíče \code{contributors\_teams} (73 řádků), poté jsou pro každého přispěvatele data čtena přímo z~pokrývajícího indexu \code{idx\_file\_activity\_repo\_contrib\_ts} s~\code{Heap Fetches: 0}. Zrychlení je způsobeno především eliminací paralelizace a~přeskočením desítek tisíc nesouvisejících řádků.

\begin{longlisting}
  \caption{S3: EXPLAIN ANALYZE s~indexy (Execution Time: 41,0 ms)}
  \label{lst:appendix:explain:s3:idx}
  \inputminted[fontsize=\footnotesize]{text}{./code/sql/explain_q3_idx.txt}
\end{longlisting}

% ── S4 ─────────────────────────────────────────────────────────────────────
\subsection*{S4 --- bez indexů}

Bez indexů jsou v~plánu dva sekvenční průchody: \code{Seq Scan} nad \code{issue\_labels\_history} \\ a~\code{Parallel Seq Scan} nad \code{issue\_event\_history} s~filtrací \code{is\_pr} a~příslušného repozitáře. Řazení mezivýsledků pro okenní funkce tabulky \code{issue\_labels\_history} probíhá operací \\ \code{Sort (external merge Disk: 1\,664\,kB)}, protože data se nevejdou do paměti. Odhadované náklady jsou 24\,958 a~doba vykonávání 111,4\,ms.

\begin{longlisting}
  \caption{S4: EXPLAIN ANALYZE bez indexů (Execution Time: 111,4 ms)}
  \label{lst:appendix:explain:s4:no_idx}
  \inputminted[fontsize=\scriptsize]{text}{./code/sql/explain_q4_no_idx.txt}
\end{longlisting}

\subsection*{S4 --- s~indexy}

Oproti variantě bez indexů jsou hlavní změny: \code{Seq Scan} na \code{issue\_labels\_history} nahrazen \code{Index Only Scan using idx\_pr\_labels\_repo\_label\_issue\_ts} s~\code{Incremental Sort} (quick-\\sort v~paměti místo external merge na disku). \code{Parallel Seq Scan} na \code{issue\_event\_history} nahrazen \code{Bitmap Heap Scan} za použití indexu \code{idx\_issue\_event\_history\_repo\_evt\_pr\_only}, čímž odpadá paralelní průchod celou tabulkou. Dominantní náklady zůstávají v~\code{Merge Join} pro výpočet průniku intervalů, který není přímo indexovatelný.

\begin{longlisting}
  \caption{S4: EXPLAIN ANALYZE s~indexy (Execution Time: 74,6 ms)}
  \label{lst:appendix:explain:s4:idx}
  \inputminted[fontsize=\scriptsize]{text}{./code/sql/explain_q4_idx.txt}
\end{longlisting}
